
\providecommand{\ndensity}{n^{\prime\prime\prime}}
\providecommand{\sfission}{s_f^{\prime\prime\prime}}
\providecommand{\dd}{\,\mathrm{d}}
\providecommand{\SourceEq}[1]{\par\nobreak\hfill{\sffamily\scriptsize\color{KSUGray}#1}\par}

\setlength{\abovedisplayskip}{4pt plus 1pt minus 1pt}
\setlength{\belowdisplayskip}{4pt plus 1pt minus 1pt}
\setlength{\abovedisplayshortskip}{3pt plus 1pt}
\setlength{\belowdisplayshortskip}{3pt plus 1pt}
\setlength{\jot}{2pt}
\tcbset{handout base/.append style={before skip=2.5pt,after skip=2.5pt}}

\HandoutSetup
  {NE 630}
  {09}
  {Fast and Epithermal Spectra}
  {Textbook \S3.4}

\begin{document}
\MakeHandoutTitle

\textcolor{KSUDeepPurple}{\textbf{Primary objective.}}
Approximate the flux spectrum $\phi(E)$ of fast and epithermal neutrons
subject to elastic collisions without absorption.
\par\vspace{5pt}

\begin{handoutcolumns}
\HandoutSection{1}{Review: why do we need a spectrum?}
Last time, $\eta(E)$ was a first step toward $k$: an energy average
would give a single value approximately bounding $k$ from above.
Before that, we developed the energy-transfer distribution $P(E\!\to\!E')$.

\begin{checkpoint}[The missing weighting function]
The distribution of neutron energies comes from competition between
\textbf{scattering and absorption}. What weighting function should we
use when averaging over energy?
\RuledLines{1}
\end{checkpoint}

\HandoutSection{2}{Neutron density, flux, and reaction rates}
All particle counts, distances, and interaction counts below are
\textbf{expected} quantities.

\begin{fixedbox}{Density and flux distributions}
$\ndensity(E)$ is the neutron density per unit energy, with units
$\mathrm{cm}^{-3}\,\mathrm{eV}^{-1}$. The \textbf{neutron flux} is
$$
\phi(E)=\ndensity(E)\,v(E),\qquad
v(E)=\sqrt{\frac{2E}{m}}.
$$
Here $m$ is the neutron mass and $E=\tfrac12 mv^2$.
The units of $\phi(E)$ are
$\mathrm{cm}^{-2}\,\mathrm{s}^{-1}\,\mathrm{eV}^{-1}$.
\end{fixedbox}

\textbf{Path-length interpretation.}
$\phi(E)\dd E$ is the total distance traveled in one second by all
neutrons with energies between $E$ and $E+\mathrm{d}E$ in a
$1\unit{cm^3}$ volume. This interpretation is especially useful for
\textbf{Monte Carlo simulation}.

\begin{fixedbox}{Differential and integrated reaction rates}
$$
\phi(E)\,\Sigma_x(E)\dd E
$$
is the number of interactions of type $x$, per second per
$\mathrm{cm}^3$, for neutrons in the same energy interval.
Integrating over energy gives
$$
R_x=\int_0^\infty\Sigma_x(E)\,\phi(E)\dd E,
\qquad [R_x]=\mathrm{cm}^{-3}\,\mathrm{s}^{-1}.
$$
\end{fixedbox}

\begin{checkpoint}[Units and meaning]
What distinguishes $\phi(E)$ from $\phi(E)\dd E$? What does multiplying
by $\Sigma_x(E)$ count?
\RuledLines{2}
\end{checkpoint}

\switchcolumn
\RaggedRight
\HandoutSection{3}{Balance at a neutron energy $E$}
The \textbf{total interaction rate density} is
$$
\Sigma_t(E)\,\phi(E)
\quad [\mathrm{cm}^{-3}\,\mathrm{s}^{-1}\,\mathrm{eV}^{-1}].
$$
Each interaction removes a neutron from energy $E$. To balance that
loss, neutrons must scatter into $E$, be born at $E$ by fission, or
arrive from an external source.

\begin{fixedbox}{The slowing-down (or spectrum) equation}
\begin{align*}
\Sigma_t(E)\,\phi(E)
 &=\int_0^\infty P(E'\!\to\!E)\,\Sigma_s(E')\phi(E')\dd E'\\
 &\quad+\chi(E)\int_0^\infty\bar\nu(E')\Sigma_f(E')\phi(E')\dd E'\\
 &\quad+s_{\mathrm{ext}}(E).
\end{align*}
\SourceEq{Like FNRP 3.15}
\end{fixedbox}

The right-hand terms describe \textbf{in-scattering}, \textbf{fission
births}, and an \textbf{external source}, respectively. The
energy-transfer cross section is
$$
\Sigma_s(E'\!\to\!E)=P(E'\!\to\!E)\,\Sigma_s(E').
$$
This balance follows from the energy-dependent diffusion equation
when space and time are eliminated.

\HandoutSection{4}{Three energy ranges}
\begin{fixedbox}{Approximations by energy range}
\renewcommand{\arraystretch}{1.13}
\setlength{\tabcolsep}{3pt}
\begin{tabularx}{\linewidth}{@{}l >{\raggedright\arraybackslash}X@{}}
\textbf{Fast} & $0.1$--$10\unit{MeV}$; $\chi(E)>0$.\\
 & $\Sigma_s(E'\!\to\!E)=0$ for $E'<E$.\\[3pt]
\textbf{Epithermal} & $1\unit{eV}$--$0.1\unit{MeV}$; $\chi(E)=0$.\\
 & $\Sigma_s(E'\!\to\!E)=0$ for $E'<E$.\\[3pt]
\textbf{Thermal} & $10^{-3}$--$1\unit{eV}$; $\chi(E)=0$.\\
 & $\Sigma_s(E'\!\to\!E)\ne0$ can occur for $E'<E$.
\end{tabularx}
\end{fixedbox}

\textbf{Why can thermal neutrons gain energy?}
Their kinetic energies can be comparable to those of surrounding
nuclei, so a collision can transfer energy \emph{to} the neutron.

\textbf{Why neglect low-energy fission births?}
Most fission neutrons have energies around $1$--$2\unit{MeV}$;
very few are born below $0.1\unit{MeV}$.

\begin{checkpoint}[Read the direction of transfer]
In $P(E'\!\to\!E)$, which energy is before the collision? Why does
$E'<E$ describe an energy gain?
\RuledLines{1}
\end{checkpoint}
\end{handoutcolumns}

\newpage

\begin{handoutcolumns}
\HandoutSection{5}{A first approximation to the fast spectrum}
The \textbf{uncollided flux} in the fast range provides a first
approximation to the fast spectrum. We label this approximation
$\phi_{\mathrm{fast}}(E)$ here.

\begin{fixedbox}{Uncollided fast-spectrum approximation}
$$
\phi_{\mathrm{fast}}(E)\simeq
\frac{\chi(E)}{\Sigma_t(E)}\,\sfission.
$$
\SourceEq{FNRP 3.16}
The fission-source strength is
$$
\sfission=\int_0^\infty\nu(E')\Sigma_f(E')\phi(E')\dd E'.
$$
\end{fixedbox}

\begin{checkpoint}[The flux depends on the flux!]
Where does $\phi$ appear on the right-hand side? Which factors set
the spectral shape, and which quantity sets its amplitude?
\RuledLines{2}
\end{checkpoint}

\textbf{Scattering changes the fast spectrum.}
Any light nuclei cause energy degradation through elastic scattering.
Inelastic scattering with heavy nuclei is also very important.

\begin{boardbox}{Board sketch: inelastic-scattering energy losses}
\begin{minipage}[t][0.65in][t]{\linewidth}
{\footnotesize\itshape Sketch the inelastic-scattering energy changes, $\Delta E$.}
\end{minipage}
\end{boardbox}

\HandoutSection{6}{Epithermal slowing down: assumptions}
Return to the spectrum balance and assume:
\begin{enumerate}[leftmargin=1.55em,itemsep=1pt,topsep=2pt]
  \item $E<0.1\unit{MeV}$: no direct fission births at $E$.
  \item No absorption.
  \item No external source term in this energy range.
\end{enumerate}
Absorption is deferred to the next lesson.

\begin{fixedbox}{Scattering-only balance}
$$
\Sigma_s(E)\phi(E)=
\int_0^\infty P(E'\!\to\!E)\Sigma_s(E')\phi(E')\dd E'.
$$
\SourceEq{FNRP 3.21}
\end{fixedbox}

\begin{boardbox}{OpenMC example / observations}
\RuledLines{3}
\end{boardbox}

\switchcolumn
\RaggedRight
\HandoutSection{7}{The energy-transfer kernel and its limits}
For a single moderator species of mass number $A$, recall
$$
\alpha=\frac{(A-1)^2}{(A+1)^2}.
$$
\begin{fixedbox}{Elastic-scattering energy distribution}
$$
P(E'\!\to\!E)=
\begin{cases}
\dfrac{1}{(1-\alpha)E'},&\alpha E'<E<E',\\[4pt]
0,&\text{otherwise}.
\end{cases}
$$
\end{fixedbox}
\textbf{Note the switched prime!} $E'$ is now the incident energy;
$E$ is the outgoing energy at which we write the balance.

\begin{checkpoint}[Exercise: determine the integration limits]
With $E$ fixed, solve $\alpha E'<E<E'$ for the incident energies
$E'$ that can contribute. For $\alpha>0$,
$$
\MathBlank[0.45in]{}<E'<\MathBlank[0.55in]{}.
$$
\RuledLines{1}
\end{checkpoint}

The resulting reduced balance is
$$
\Sigma_s(E)\phi(E)=
\int_E^{E/\alpha}\frac{\Sigma_s(E')\phi(E')}{(1-\alpha)E'}\dd E'.
$$

\HandoutSection{8}{Spectrum and slowing-down density}
\begin{fixedbox}{Forms to examine}
Candidate shape to verify:
$$
\phi(E)=\frac{C}{E}.
$$
In terms of the slowing-down density:
$$
\boxed{\displaystyle\phi(E)=\frac{q}{\xi\,\Sigma_s(E)\,E}}.
$$
Here $q$ is the \textbf{slowing-down density} and $\xi$ is the
slowing-down decrement. In the absence of absorption,
$$
q=\sfission.
$$
\end{fixedbox}

\begin{checkpoint}[Verification and consistency check]
Substitute the proposed $C/E$ form into the reduced balance.
Compare it with $q/[\xi\Sigma_s(E)E]$: what must be true of
$\Sigma_s(E)$ for both forms to describe the same flux with constant
$C$, $q$, and $\xi$?
\RuledLines{3}
\end{checkpoint}
\end{handoutcolumns}

\end{document}
